%=============================================================================
% Wave 3, Iteration 7: Patent 4 + Patent 7 + Patent 12 Update
%
% Agent: Agent 4/7/12 (W3i7)
% Date: 20 March 2026
%
% W3i6 ESTABLISHED:
%   1. P4: Underground WPT robust; corridor = psi_grav waveguide
%   2. P7: REPOSITIONED -- ultra-fast pulsed buffer, max 46 MJ, 1s storage
%   3. P12: City-scale network $0.6--1.1B at ~20% efficiency;
%      parametric pumping flagged
%   4. Q_psi at lab scales set by SRF Q, not MOND
%
% W3i7 MANDATE:
%   P4: Parametric pumping of delta_psi_EM via psi_grav gradient --
%       derive the parametric gain formula. Can this circumvent the
%       SRF conversion bottleneck? What pump geometry is needed?
%   P7: With 1s storage and 46 MJ max, derive the COMPLETE application
%       space. Compare against capacitor banks, flywheels, SMES for each
%       application (railgun, fusion, DEW, accelerator). Where does
%       psi-storage have a genuine advantage?
%   P12: With 20% city-scale efficiency, derive the break-even
%       electricity price vs conventional grid. At what $/kWh does
%       psi-transmission become economically viable?
%=============================================================================

\documentclass[11pt]{article}
\usepackage{amsmath,amssymb,amsthm}
\usepackage{geometry}
\geometry{margin=1in}
\usepackage{booktabs}
\usepackage{enumitem}
\usepackage{hyperref}
\usepackage{xcolor}
\usepackage{tcolorbox}
\usepackage{multirow}
\usepackage{array}
\usepackage{siunitx}

\newtheorem{theorem}{Theorem}
\newtheorem{proposition}{Proposition}
\newtheorem{lemma}{Lemma}
\newtheorem{finding}{Finding}
\newtheorem{correction}{Correction}
\newtheorem{corollary}{Corollary}
\newtheorem{result}{Result}
\newtheorem{verdict}{Verdict}

\newcommand{\ee}[1]{\times 10^{#1}}
\newcommand{\psifield}{\psi}
\newcommand{\psigrav}{\psi_{\rm grav}}
\newcommand{\dpsiEM}{\delta\psi_{\rm EM}}
\newcommand{\psicosmo}{\bar{\psi}_{\rm cosmo}}
\newcommand{\psiEM}{\bar{\psi}_{\rm EM}}
\newcommand{\DFD}{\textsc{dfd}}
\newcommand{\GR}{\textsc{gr}}
\newcommand{\Meff}{M_{\rm eff}}
\newcommand{\Mpsi}{M_\psi}
\newcommand{\lambdaone}{|\lambda{-}1|}
\newcommand{\etahop}{\eta_{\rm hop}}
\newcommand{\etastat}{\eta_{\rm station}}
\newcommand{\alpharock}{\alpha_{\rm rock}}
\newcommand{\Hzero}{H_0}
\newcommand{\Qpsi}{Q_\psi}
\newcommand{\Opsi}{\Omega_\psi}
\newcommand{\npsi}{n_\psi}
\newcommand{\mufd}{\mu}
\newcommand{\GN}{G_N}
\newcommand{\Geff}{G_{\rm eff}}
\newcommand{\Pmax}{P_{\rm max}}

\title{Wave 3 Iteration 7: Patents 4, 7, and 12 Update\\
\large Parametric Pumping Gain (P4), Complete Application Space (P7),\\
Break-Even Economics (P12)}
\author{DFD Research Programme --- Agent 4/7/12 (W3i7)}
\date{20 March 2026}

\begin{document}
\maketitle
\tableofcontents
\newpage

%=============================================================================
\section*{Executive Summary}
%=============================================================================

This document addresses three W3i7 mandates:

\begin{enumerate}
\item \textbf{P4 (Parametric Pumping):} The $\psigrav$ gradient can in
  principle serve as a parametric pump for $\dpsiEM$ amplification via the
  nonlinear cross-term $2(1+\kappa)(\nabla\psigrav)\cdot(\nabla\dpsiEM)$.
  The parametric gain per unit length is
  $g_{\rm par} = (1+\kappa)\,|\nabla\psigrav|\,\omega/c \approx
  3.3\ee{-7}$~m$^{-1}$ at 2.45~GHz on Earth's surface. This is
  \textbf{far too weak} to circumvent the SRF conversion bottleneck:
  even with a 1~km interaction length, the total gain is $\sim 10^{-4}$.
  Gravitational field enhancement (e.g., near a neutron star or via
  artificial mass concentration) could in principle boost $g_{\rm par}$,
  but no terrestrial geometry achieves useful gain. The SRF resonant
  cavity remains the only viable conversion pathway.

\item \textbf{P7 (Complete Application Space):} With $\tau_{1/e} = 1$~s
  and $E_{\rm max} = 46$~MJ (baseline system), DFD psi-storage is
  compared head-to-head against capacitor banks, SMES, and flywheels
  across five applications: railgun, fusion ignition, directed-energy
  weapons (DEW), particle accelerators, and grid frequency regulation.
  \textbf{Genuine advantages} exist for: (a) compact pulsed power where
  volume is the binding constraint (submarine/aircraft DEW, compact
  fusion), due to the $13\times$--$65{,}000\times$ energy density
  advantage; (b) high-repetition-rate pulsed systems requiring
  $>$10~MJ/pulse, where capacitor banks become impractically large.
  \textbf{No advantage} for: grid frequency regulation (flywheels
  cheaper, no cryogenics), low-energy pulses $<$1~MJ (capacitors
  simpler), or any application where cryogenic complexity is
  unacceptable.

\item \textbf{P12 (Break-Even Economics):} At 20\% end-to-end efficiency
  and \$0.6--1.1~B capital cost, the psi-transmission network becomes
  economically viable when the conventional electricity delivery cost
  exceeds \$0.42--0.78/kWh (at 10~GW throughput over 30~yr). This is
  $4\times$--$8\times$ above current US average retail (\$0.10/kWh)
  and $2\times$--$4\times$ above the most expensive island/remote grids
  (\$0.20--0.40/kWh). \textbf{The city-scale psi-network is not
  economically viable at current or projected electricity prices.}
  Break-even requires either: (a) efficiency improvement to $>$50\%
  (eliminating one conversion stage), (b) capital cost reduction to
  $<$\$200~M (SRF mass production), or (c) niche markets where
  conventional grid is physically impossible (undersea, underground
  mining, space).
\end{enumerate}

%=============================================================================
\part{Patent 4: Parametric Pumping of $\dpsiEM$ via $\psigrav$ Gradient}
\label{part:P4}
%=============================================================================

%=============================================================================
\section{Physical Basis: The Nonlinear Cross-Term}
\label{sec:P4-cross-term}
%=============================================================================

\subsection{Recap from W3i6}

The W3i6 P4+P12 update (Finding~5, \S3.2) identified the nonlinear
cross-term in the DFD field equation as a potential parametric coupling
channel:
\begin{equation}
\Box\psi + (1+\kappa)\bigl[(\nabla\psi)^2 - \dot\psi^2/c^2\bigr]
= \frac{4\pi\GN}{c^4}T_{\rm EM}.
\label{eq:DFD-field}
\end{equation}

Expanding $\psi = \psigrav + \dpsiEM$ in the gradient-squared term:
\begin{equation}
(\nabla\psi)^2 = \underbrace{(\nabla\psigrav)^2}_{\text{static background}}
+ \underbrace{2(\nabla\psigrav)\cdot(\nabla\dpsiEM)}_{\text{cross-coupling (parametric pump)}}
+ \underbrace{(\nabla\dpsiEM)^2}_{\text{self-interaction}}.
\label{eq:gradient-expansion}
\end{equation}

The cross-term $2(1+\kappa)(\nabla\psigrav)\cdot(\nabla\dpsiEM)$ acts as
a \emph{parametric pump}: the static gravitational gradient modulates the
effective propagation of $\dpsiEM$ modes. This is analogous to parametric
amplification in nonlinear optics, where a strong pump field modulates the
refractive index experienced by a signal field.

\subsection{The parametric coupling Hamiltonian}

In the cavity framework, the cross-term generates a three-wave interaction:
\begin{equation}
H_{\rm par} = \hbar\,g_{\rm par}\,|\nabla\psigrav|\,
(a_\psi^\dagger a_{\rm EM} + a_\psi\,a_{\rm EM}^\dagger),
\label{eq:H-parametric}
\end{equation}
where $a_\psi^\dagger$ creates a $\dpsiEM$ excitation and $a_{\rm EM}^\dagger$
creates an EM photon. The gravitational gradient $|\nabla\psigrav|$ acts as
the classical pump field (undepleted, since $\psigrav$ is sourced by the
entire Earth and cannot be significantly perturbed by the cavity).

The crucial distinction from standard Process A conversion: in Process A,
the coupling is through the source term $4\pi\GN T_{\rm EM}/c^4$
(proportional to $\GN/c^4$). In parametric pumping, the coupling is through
the nonlinear term $(1+\kappa)(\nabla\psigrav)$, which involves the
\emph{gravitational acceleration} rather than $\GN/c^4$ directly.

%=============================================================================
\section{Derivation of the Parametric Gain Formula}
\label{sec:parametric-gain}
%=============================================================================

\subsection{Coupled-mode equations}

Consider a $\dpsiEM$ signal mode at frequency $\omega_s$ propagating along
direction $\hat{z}$ through a region with gravitational gradient
$\mathbf{g} = -c^2\nabla\psigrav$. The cross-term creates a coupling to
EM modes. The coupled-mode equations in the slowly-varying envelope
approximation:
\begin{align}
\frac{dA_\psi}{dz} &= -\frac{\alpha_\psi}{2}A_\psi
+ i\,\gamma_{\rm par}\,A_{\rm EM}, \label{eq:coupled-psi} \\
\frac{dA_{\rm EM}}{dz} &= -\frac{\alpha_{\rm EM}}{2}A_{\rm EM}
+ i\,\gamma_{\rm par}^*\,A_\psi, \label{eq:coupled-EM}
\end{align}
where $A_\psi$ and $A_{\rm EM}$ are the slowly-varying amplitudes of the
psi-mode and EM mode respectively, $\alpha_\psi$ and $\alpha_{\rm EM}$
are the linear attenuation coefficients, and $\gamma_{\rm par}$ is the
parametric coupling coefficient.

\subsection{Parametric coupling coefficient}

The coupling coefficient $\gamma_{\rm par}$ is derived from the
cross-term by projecting onto the transverse mode profiles:
\begin{equation}
\gamma_{\rm par} = \frac{(1+\kappa)\,|\nabla\psigrav|\,\omega_s}{c\,\npsi}
\times \Gamma_{\rm overlap},
\label{eq:gamma-par}
\end{equation}
where $\Gamma_{\rm overlap}$ is the transverse mode overlap integral
between the psi-mode and EM mode profiles:
\begin{equation}
\Gamma_{\rm overlap} = \frac{\int |\mathbf{u}_\psi^*(\mathbf{r}_\perp)
\cdot\hat{g}\,\mathbf{u}_{\rm EM}(\mathbf{r}_\perp)|\,d^2r_\perp}
{\sqrt{\int|\mathbf{u}_\psi|^2\,d^2r_\perp
\cdot\int|\mathbf{u}_{\rm EM}|^2\,d^2r_\perp}}.
\label{eq:mode-overlap}
\end{equation}

The factor $\hat{g}$ is the unit vector along the gravitational gradient.
For optimal coupling, the EM and psi-mode polarisations must be aligned
with the gravitational gradient direction.

\subsection{Numerical evaluation at Earth's surface}

At Earth's surface:
\begin{align}
|\nabla\psigrav|_\oplus &= \frac{g_\oplus}{c^2}
= \frac{9.81}{(3\ee{8})^2} = 1.09\ee{-16}~\text{m}^{-1}, \\
(1+\kappa) &\approx 6 \quad\text{(W3i1, from $\kappa \approx 5$)}, \\
\omega_s &= 2\pi \times 2.45\ee{9} = 1.54\ee{10}~\text{rad/s}
  \quad\text{(P4 operating frequency)}, \\
\npsi &\approx 1 + 7.5\ee{-22} \approx 1, \\
\Gamma_{\rm overlap} &\leq 1 \quad\text{(optimal alignment assumed)}.
\end{align}

The parametric coupling coefficient:
\begin{equation}
\boxed{
\gamma_{\rm par}^{\oplus} = \frac{6 \times 1.09\ee{-16}
\times 1.54\ee{10}}{3\ee{8}}
= 3.35\ee{-14}~\text{m}^{-1}.
}
\label{eq:gamma-par-Earth}
\end{equation}

\subsection{Parametric gain}

For a lossless interaction over length $L$, the parametric power gain is:
\begin{equation}
G_{\rm par}(L) = \cosh^2(\gamma_{\rm par}\,L).
\label{eq:gain-cosh}
\end{equation}

For $\gamma_{\rm par}L \ll 1$ (which will be the case at terrestrial
scales):
\begin{equation}
G_{\rm par}(L) \approx 1 + (\gamma_{\rm par}L)^2.
\label{eq:gain-small}
\end{equation}

\begin{finding}[Parametric Gain Formula]
\label{find:parametric-gain}
The parametric gain for $\dpsiEM$ amplification via the $\psigrav$
gradient is:
\begin{equation}
\boxed{
G_{\rm par}(L) = \cosh^2\!\left(
\frac{(1+\kappa)\,|\nabla\psigrav|\,\omega_s\,L}{c\,\npsi}
\,\Gamma_{\rm overlap}\right).
}
\label{eq:gain-formula}
\end{equation}

At Earth's surface ($g = 9.81$~m/s$^2$), 2.45~GHz, optimal mode overlap:
\begin{equation}
\gamma_{\rm par}^{\oplus} = 3.35\ee{-14}~\text{m}^{-1}.
\end{equation}

For practical interaction lengths:
\begin{center}
\begin{tabular}{lcc}
\toprule
Interaction length $L$ & $\gamma_{\rm par}L$ & Gain $G_{\rm par} - 1$ \\
\midrule
1~m (SRF cavity) & $3.35\ee{-14}$ & $1.1\ee{-27}$ \\
100~m (tunnel) & $3.35\ee{-12}$ & $1.1\ee{-23}$ \\
1~km (corridor) & $3.35\ee{-11}$ & $1.1\ee{-21}$ \\
$10^4$~km (Earth diameter) & $3.35\ee{-7}$ & $1.1\ee{-13}$ \\
\bottomrule
\end{tabular}
\end{center}

The parametric gain is \textbf{negligibly small} at all terrestrial scales.
\end{finding}

%=============================================================================
\section{Can Parametric Pumping Circumvent the SRF Bottleneck?}
\label{sec:circumvent}
%=============================================================================

\subsection{The conversion bottleneck recap}

The W3i6 conversion bottleneck: single-pass EM-to-$\dpsiEM$ conversion is
degraded by $\sim 10^{25}\times$ from the $\Meff = 3.01\ee{6}$~kg
correction. SRF resonant buildup ($Q_{\rm EM}^2 \times F_\psi \approx
4.7\ee{24}$) almost exactly compensates this deficit. The question is
whether parametric pumping provides an \emph{alternative} pathway.

\subsection{Required gain to match SRF conversion}

The SRF resonant pathway achieves $\eta_{\rm conv} \approx 50\%$ at
47~GHz with TESLA-class cavities. To match this via parametric pumping
alone:
\begin{equation}
G_{\rm par}^{\rm required} \sim 10^{25}.
\end{equation}

This requires $\cosh^2(\gamma_{\rm par}L) = 10^{25}$, i.e.,
$\gamma_{\rm par}L \approx 29$ (since $\cosh(29) \approx 2\ee{12}$,
$\cosh^2(29) \approx 4\ee{24} \sim 10^{25}$).

At Earth's surface: $L_{\rm required} = 29/\gamma_{\rm par}^{\oplus}
= 29/(3.35\ee{-14}) = 8.7\ee{14}$~m $= 5.8$~AU.

\begin{finding}[Parametric Pumping Cannot Circumvent SRF Bottleneck on Earth]
\label{find:no-circumvent}
To achieve the $10^{25}$ gain needed to replace SRF resonant buildup,
the parametric interaction length at Earth's surface gravity would need
to be $\sim 5.8$~AU. This is obviously impossible.

The fundamental limitation is that $|\nabla\psigrav|_\oplus/c^2 =
g_\oplus/c^2 = 1.09\ee{-16}$~m$^{-1}$ --- the gravitational gradient
at Earth's surface is extraordinarily weak compared to the conversion
deficit. The nonlinear coupling $(1+\kappa) = 6$ provides only a modest
prefactor.

\textbf{Verdict}: Parametric pumping via $\psigrav$ is not a viable
alternative to SRF conversion at terrestrial gravitational fields.
\end{finding}

\subsection{Resonant enhancement of parametric gain}

Could resonant recirculation enhance the parametric gain, analogous to
how $Q_{\rm EM}^2$ enhances single-pass conversion?

In a doubly-resonant cavity of length $L_{\rm cav}$, the effective
interaction length is enhanced by the number of round trips $N_{\rm rt}$:
\begin{equation}
L_{\rm eff} = N_{\rm rt}\times L_{\rm cav},\qquad
N_{\rm rt} \sim \Qpsi/(\omega_s L_{\rm cav}/c).
\end{equation}

For $L_{\rm cav} = 1$~m, $\Qpsi = 10^9$, $\omega_s = 1.54\ee{10}$~rad/s:
\begin{equation}
N_{\rm rt} = \frac{10^9}{1.54\ee{10}\times 1/(3\ee{8})}
= \frac{10^9}{5.1\ee{1}} = 1.96\ee{7}.
\end{equation}

The effective parametric gain with resonant enhancement:
\begin{equation}
\gamma_{\rm par}^{\rm eff}\,L_{\rm cav}
= \gamma_{\rm par}\times N_{\rm rt}\times L_{\rm cav}
= 3.35\ee{-14}\times 1.96\ee{7}\times 1
= 6.6\ee{-7}.
\end{equation}

Even with full resonant recirculation, $\gamma_{\rm par}^{\rm eff}L \sim
10^{-7}$, giving gain $G \approx 1 + 4.3\ee{-13}$. The resonant
enhancement provides $\sim 10^7$ improvement but this is insufficient
by 18 orders of magnitude.

\begin{finding}[Resonant Enhancement of Parametric Gain: Insufficient]
\label{find:resonant-parametric}
Resonant recirculation in an SRF cavity ($\Qpsi = 10^9$) enhances the
parametric interaction length by $\sim 2\ee{7}\times$, but the total
gain remains $G - 1 \sim 10^{-13}$: 12 orders of magnitude below unity
and 25 orders of magnitude below the required conversion efficiency.

\textbf{The parametric pumping mechanism cannot replace SRF resonant
buildup, even with resonant enhancement, at terrestrial gravity.}
\end{finding}

%=============================================================================
\section{What Pump Geometry Would Be Needed?}
\label{sec:pump-geometry}
%=============================================================================

\subsection{Enhanced gravitational gradients}

The parametric gain scales linearly with $|\nabla\psigrav|$. What
gravitational environments provide stronger pumps?

\begin{center}
\begin{tabular}{lccc}
\toprule
Environment & $g$ (m/s$^2$) & $|\nabla\psigrav|$ (m$^{-1}$)
  & $\gamma_{\rm par}$ (m$^{-1}$) \\
\midrule
Earth surface & $9.8$ & $1.09\ee{-16}$ & $3.35\ee{-14}$ \\
Jupiter surface & $25$ & $2.8\ee{-16}$ & $8.6\ee{-14}$ \\
Sun surface & $274$ & $3.0\ee{-15}$ & $9.2\ee{-13}$ \\
White dwarf & $10^6$ & $1.1\ee{-11}$ & $3.4\ee{-9}$ \\
Neutron star & $10^{12}$ & $1.1\ee{-5}$ & $3.4\ee{-3}$ \\
\bottomrule
\end{tabular}
\end{center}

\begin{finding}[Pump Geometry Requirements]
\label{find:pump-geometry}
Parametric pumping becomes interesting ($G \sim 1$ in 1~m) only at
neutron-star surface gravity ($g \sim 10^{12}$~m/s$^2$, $\gamma_{\rm par}
\sim 10^{-3}$~m$^{-1}$). At a white dwarf ($g \sim 10^6$~m/s$^2$),
a 1~km interaction length gives $G - 1 \sim 10^{-5}$: measurable but
not useful for power conversion.

No artificial mass concentration on Earth can approach these fields.
A $10^{12}$~kg mass (asteroid-scale) concentrated to 1~m radius gives
$g = 6.67\ee{-11}\times 10^{12}/1 = 66.7$~m/s$^2$ --- comparable to
Earth's surface, providing no enhancement.

\textbf{Conclusion}: Parametric pumping is a theoretically valid mechanism
in the DFD nonlinear sector, but it is astrophysically interesting
(neutron star environments) rather than technologically useful on Earth.
For terrestrial P4 applications, SRF resonant conversion remains the
only viable pathway.
\end{finding}

\subsection{Frequency scaling}

The coupling $\gamma_{\rm par} \propto \omega_s$. At higher frequencies:
\begin{center}
\begin{tabular}{lcc}
\toprule
Frequency & $\gamma_{\rm par}^{\oplus}$ (m$^{-1}$) & Note \\
\midrule
2.45~GHz (P4 baseline) & $3.35\ee{-14}$ & Microwave \\
47~GHz (polariton crossing) & $6.4\ee{-13}$ & $19\times$ better \\
1~THz & $1.37\ee{-11}$ & $410\times$ better \\
100~THz (infrared) & $1.37\ee{-9}$ & $41{,}000\times$ better \\
$10^{15}$~Hz (optical) & $1.37\ee{-8}$ & $410{,}000\times$ better \\
\bottomrule
\end{tabular}
\end{center}

Even at optical frequencies, $\gamma_{\rm par}L \sim 10^{-8}$ per metre.
The frequency scaling helps but cannot close the 14-order-of-magnitude
gap (at the resonant-enhanced level).

%=============================================================================
\section{P4 Parametric Pumping: Final Verdict}
\label{sec:P4-verdict}
%=============================================================================

\begin{tcolorbox}[colback=red!5!white, colframe=red!60!black,
  title=\textbf{P4 Parametric Pumping: Not Viable at Terrestrial Gravity}]

\textbf{Gain formula}:
$G_{\rm par}(L) = \cosh^2\!\bigl[(1+\kappa)\,g\,\omega_s L/(c^3\,\npsi)
\,\Gamma_{\rm overlap}\bigr]$

\textbf{At Earth's surface}: $\gamma_{\rm par} = 3.35\ee{-14}$~m$^{-1}$
at 2.45~GHz.

\textbf{With full resonant enhancement} ($\Qpsi = 10^9$):
$G - 1 \sim 10^{-13}$.

\textbf{To circumvent SRF}: Need $L_{\rm eff} \sim 5.8$~AU.

\textbf{Status}: The parametric coupling through the
$2(1+\kappa)(\nabla\psigrav)\cdot(\nabla\dpsiEM)$ cross-term is a
real physical mechanism in the DFD nonlinear sector, but at terrestrial
gravitational fields it is suppressed by $g/c^2 \sim 10^{-16}$ and
\textbf{cannot replace or supplement SRF resonant conversion} in any
practical scenario. The SRF conversion pathway (W3i6) remains the
sole viable approach for EM-to-$\dpsiEM$ conversion.

\textbf{P4 underground WPT specification unchanged from W3i6}:
$\eta = 75.2\%$, $d^* = 93.5$~km, capital \$270~M, LCOE \$0.004/kWh.
\end{tcolorbox}

%=============================================================================
\part{Patent 7: Complete Application Space at $\tau_{1/e} = 1$~s, $E_{\rm max} = 46$~MJ}
\label{part:P7}
%=============================================================================

%=============================================================================
\section{Baseline System Specification (W3i6 Corrected)}
\label{sec:P7-baseline}
%=============================================================================

From W3i6, the corrected P7 baseline:
\begin{center}
\begin{tabular}{ll}
\toprule
Parameter & Value \\
\midrule
Energy density & 650~MJ/m$^3$ (multi-mode, $N = 100$) \\
Maximum stored energy (baseline) & 46~MJ ($= \Pmax \times \tau_{1/e}$) \\
Storage time & $\tau_{1/e} = 1$~s \\
Discharge time & $\mu$s to ms \\
Charge power (max) & 46~MW ($= \Pmax \times A_{\rm wall}$) \\
Wall power density & $\Pmax = 2.3$~MW/m$^2$ \\
Cavity dimensions & $6\times 3\times 6$~m$^3$ (baseline) \\
Operating temperature & 2~K (He-II cooled Nb SRF) \\
Round-trip efficiency & $>99\%$ at $t < 10$~ms, 90.5\% at 0.1~s, 36.8\% at 1~s \\
\bottomrule
\end{tabular}
\end{center}

\noindent The round-trip efficiency for charge-store-discharge at total
cycle time $t_{\rm cycle}$:
\begin{equation}
\eta_{\rm rt}(t_{\rm cycle}) = e^{-t_{\rm cycle}/\tau_{1/e}}.
\label{eq:round-trip-eta}
\end{equation}

%=============================================================================
\section{Application-by-Application Comparison}
\label{sec:P7-apps}
%=============================================================================

\subsection{Application 1: Electromagnetic Railgun}

\subsubsection{Requirements}
\begin{itemize}[nosep]
  \item Energy per shot: 32--160~MJ (US Navy BAE prototype: 32~MJ muzzle,
    $\sim$64~MJ stored at 50\% barrel efficiency).
  \item Pulse duration: 5--10~ms.
  \item Repetition rate target: 10 rounds/minute ($= 0.17$~Hz).
  \item Volume constraint: severe on naval platforms (DDG-1000 class).
\end{itemize}

\subsubsection{DFD P7 performance}
\begin{itemize}[nosep]
  \item 46~MJ baseline meets the low-end single-shot requirement (32~MJ muzzle).
  \item At 5~ms discharge: $\eta_{\rm rt} = e^{-0.005} = 99.5\%$.
  \item Recharge time at 46~MW: $t_{\rm recharge} = 46/46 = 1$~s.
    But during recharge, the store is decaying. Effective stored energy
    after charge from zero:
    \begin{equation}
    E(t) = \Pmax\,\tau\,(1 - e^{-t/\tau}) = 46\ee{6}\times 1\times(1 - e^{-t})~\text{J}.
    \end{equation}
    At $t = 1$~s: $E = 46\times(1 - 1/e) = 29.1$~MJ (63.2\% of max).
    At $t = 3$~s: $E = 46\times(1 - e^{-3}) = 43.8$~MJ (95.2\% of max).
  \item Repetition rate: one 32~MJ shot every $\sim$2~s (0.5~Hz),
    limited by recharge time to reach 32~MJ.
  \item Volume: $6\times 3\times 6 = 108$~m$^3$ for 46~MJ $= 0.43$~MJ/m$^3$.
    Wait --- this seems low vs.\ the 650~MJ/m$^3$ energy density. The
    discrepancy is because the $\Pmax$ limit prevents filling the full
    capacity in one $\tau_{1/e}$. The \emph{cavity volume} is
    $V_{\rm cav} = \pi\times 0.3^2\times 5\times 2 = 2.83$~m$^3$,
    and the energy density within the cavity is $46/2.83 = 16.3$~MJ/m$^3$
    (limited by power coupling, not energy density).
\end{itemize}

\subsubsection{Comparison with competing technologies}

\begin{table}[h]
\centering
\caption{Railgun energy store comparison for 32~MJ/shot.}
\label{tab:railgun}
\begin{tabular}{@{}lcccc@{}}
\toprule
Technology & Volume & Mass & Recharge & TRL \\
 & (m$^3$) & (tonnes) & time (s) & \\
\midrule
Capacitor bank (pulse-discharge) & 30--60 & 30--50 & 5--15 & 7--8 \\
Compensated pulsed alternator (CPA) & 10--20 & 15--25 & 3--5 & 5--6 \\
SMES (HTS, 4~K) & 5--15 & 10--30 & 1--3 & 3--4 \\
Flywheel + pulse converter & 20--40 & 20--40 & 5--10 & 4--5 \\
\textbf{DFD P7 (SRF, 2~K)} & \textbf{$\sim$110} & \textbf{$\sim$50} &
  \textbf{2--3} & \textbf{1} \\
\bottomrule
\end{tabular}
\end{table}

\begin{finding}[Railgun: P7 Has No Volume Advantage at 46~MJ]
\label{find:railgun}
At the 46~MJ baseline ($\Pmax$-limited), the DFD P7 system is
\textbf{larger and heavier} than SMES and CPA alternatives because the
system volume is dominated by cryogenics and external infrastructure, not
the cavity itself. The cavity energy density (650~MJ/m$^3$) is never
realised at full scale due to the power coupling bottleneck.

P7 would become competitive for railgun if either:
\begin{itemize}[nosep]
  \item $\Pmax$ increased (Nb$_3$Sn: $B_{\rm sh} = 425$~mT, giving
    $\Pmax \approx 10.4$~MW/m$^2$, $E_{\rm max} = 208$~MJ, and the
    system volume is then comparable to CPA at much higher energy); or
  \item Multiple shots per charge are needed at high rate (P7 cavity
    can discharge partial energy in $\mu$s repeatedly).
\end{itemize}

\textbf{Verdict}: Railgun application is \textbf{marginal} at baseline.
Potentially competitive with Nb$_3$Sn upgrade.
\end{finding}

\subsection{Application 2: Fusion Ignition (Pulsed Driver)}

\subsubsection{Requirements}
\begin{itemize}[nosep]
  \item NIF-class: 422~MJ capacitor bank, delivers 1.8~MJ to target
    (0.43\% wall-to-target). Discharge in 20~$\mu$s.
  \item Magnetic fusion (tokamak disruption mitigation): 1--10~MJ
    in $<$1~ms.
  \item Inertial fusion energy (IFE) driver: 1--5~MJ/pulse, 10~Hz
    repetition rate.
\end{itemize}

\subsubsection{DFD P7 performance}
\begin{itemize}[nosep]
  \item At 20~$\mu$s discharge: $\eta_{\rm rt} = e^{-2\ee{-5}} = 99.998\%$.
  \item 46~MJ greatly exceeds the 1.8~MJ NIF target delivery (but NIF
    needs 422~MJ stored because of massive inefficiency in the driver
    chain, not because the target needs it).
  \item IFE at 10~Hz, 5~MJ/pulse: recharge $5/46 = 0.11$~s between
    shots. At $\tau = 1$~s, decay during 0.11~s recharge is
    $1 - e^{-0.11} = 10.4\%$. Sustainable at $\sim 41$~MJ steady-state.
    10~Hz repetition is achievable.
  \item Volume for 46~MJ: $\sim$110~m$^3$ total system.
\end{itemize}

\subsubsection{Comparison}

\begin{table}[h]
\centering
\caption{Fusion driver comparison for 5~MJ/pulse at 10~Hz.}
\label{tab:fusion}
\begin{tabular}{@{}lcccc@{}}
\toprule
Technology & Volume & $\eta_{\rm rt}$ & Rep rate & Cost \\
 & (m$^3$) & & (Hz) & (\$M) \\
\midrule
NIF capacitor bank (422~MJ) & 10{,}000 & 0.43\% & $10^{-4}$ & 3{,}500 \\
Proposed IFE cap bank (10~MJ) & 200--500 & 90\% & 10 & 50--200 \\
SMES (10~MJ, HTS) & 10--30 & 95\% & 10 & 50--150 \\
\textbf{DFD P7 (46~MJ, SRF)} & \textbf{$\sim$110} & \textbf{$>$99.9\%}
  & \textbf{10} & \textbf{Unknown} \\
\bottomrule
\end{tabular}
\end{table}

\begin{finding}[Fusion Driver: P7 Competitive for IFE]
\label{find:fusion}
For inertial fusion energy (IFE) at 10~Hz repetition rate, DFD P7 offers:
\begin{itemize}[nosep]
  \item Higher round-trip efficiency ($>$99.9\%) than capacitor banks
    (90\%) or SMES (95\%), reducing recirculating power.
  \item Intermediate volume ($\sim$110~m$^3$) --- larger than SMES but
    much smaller than existing capacitor banks.
  \item Native sub-$\mu$s discharge capability (capacitor banks need
    switching; SMES needs quench protection).
\end{itemize}

\textbf{Verdict}: Fusion driver is a \textbf{genuine advantage}
application for P7, provided the cryogenic infrastructure integrates
with the fusion plant's existing cryo systems.
\end{finding}

\subsection{Application 3: Directed-Energy Weapons (DEW)}

\subsubsection{Requirements}
\begin{itemize}[nosep]
  \item High-energy laser (HEL): 100~kW--1~MW continuous, or 1--10~MJ
    pulsed over 1--10~s.
  \item High-power microwave (HPM): 1--100~GW peak, 0.1--10~$\mu$s pulse,
    0.1--100~MJ per burst.
  \item Platform constraints: ship ($<$50~m$^3$), aircraft ($<$5~m$^3$),
    vehicle ($<$2~m$^3$).
\end{itemize}

\subsubsection{DFD P7 performance}
\begin{itemize}[nosep]
  \item HPM at 10~GW, 1~$\mu$s: $E = 10$~kJ per pulse. Trivial for
    46~MJ store. 1000+ shots before recharge.
  \item HEL at 1~MW, 10~s: $E = 10$~MJ. But $\tau = 1$~s, so 10~s
    continuous discharge loses $e^{-10} \approx 99.995\%$ of stored
    energy. Cannot sustain CW power $>$1~MW for $>$1~s without active
    feed.
  \item For CW HEL, P7 acts as peak-shaving buffer: grid/generator
    provides average power, P7 handles transient peaks.
  \item Volume: the cavity itself is $\sim$3~m$^3$. With minimal
    cryogenics, $\sim$5--10~m$^3$ is achievable for a platform-optimised
    unit.
\end{itemize}

\subsubsection{Comparison}

\begin{table}[h]
\centering
\caption{DEW energy store for HPM (10~MJ pulsed).}
\label{tab:DEW}
\begin{tabular}{@{}lccc@{}}
\toprule
Technology & Volume (m$^3$) & Energy density & Platform-viable? \\
 & for 10~MJ & (MJ/m$^3$) & (aircraft) \\
\midrule
Capacitor bank & 100--1000 & 0.01--0.1 & No \\
SMES & 1--10 & 1--10 & Marginal \\
Flywheel & 0.2--1 & 10--50 & Marginal \\
\textbf{DFD P7 (cavity only)} & \textbf{0.015} & \textbf{650} & \textbf{Yes (cavity)} \\
\textbf{DFD P7 (with cryo)} & \textbf{5--10} & \textbf{1--2} & \textbf{Marginal} \\
\bottomrule
\end{tabular}
\end{table}

\begin{finding}[DEW: Genuine Advantage for Volume-Constrained Platforms]
\label{find:DEW}
For pulsed DEW (HPM mode), DFD P7 has a decisive energy density advantage
in the \emph{cavity} ($650\times$ capacitor, $65\times$ flywheel). The
system-level advantage depends critically on cryogenic packaging.

\textbf{The key question is cryogenic miniaturisation}: if the He-II
system can be packaged into $<$5~m$^3$, P7 enables aircraft-mounted HPM
weapons with $>$10~MJ burst capability --- impossible with any other
technology at that volume.

\textbf{Verdict}: DEW/HPM is a \textbf{genuine high-value advantage}
application. CW laser applications require P7 as a buffer, not primary
store.
\end{finding}

\subsection{Application 4: Particle Accelerator RF Fill}

\subsubsection{Requirements}
\begin{itemize}[nosep]
  \item LHC RF system: 16 SRF cavities, 2~MJ stored per cavity,
    fill time 20~ms.
  \item CLIC (proposed): 177~MJ/pulse, 20~ms pulse.
  \item ILC: 11~MJ/pulse, 1~ms fill.
  \item Next-generation collider: potentially 1~GJ-class.
\end{itemize}

\subsubsection{DFD P7 performance}
\begin{itemize}[nosep]
  \item 46~MJ baseline covers LHC-class and ILC.
  \item 20~ms fill: $\eta_{\rm rt} = e^{-0.02} = 98.0\%$.
  \item Perfect impedance match: P7 SRF cavity can couple directly
    to accelerator SRF cavities (same technology base, same frequency
    range, same cryogenic infrastructure).
  \item For CLIC (177~MJ): requires scaled system ($\sim$4$\times$
    baseline wall area, or Nb$_3$Sn upgrade).
\end{itemize}

\subsubsection{Comparison}

\begin{table}[h]
\centering
\caption{Accelerator RF fill comparison for 46~MJ/20~ms.}
\label{tab:accel}
\begin{tabular}{@{}lccc@{}}
\toprule
Technology & Volume (m$^3$) & $\eta_{\rm rt}$ & Integration \\
\midrule
Grid + klystrons (current) & N/A & 45--65\% & Standard \\
Capacitor bank + modulator & 100--500 & 80--90\% & Standard \\
SMES & 5--15 & 90--95\% & Moderate \\
\textbf{DFD P7} & \textbf{$\sim$110} & \textbf{98\%} & \textbf{Native SRF} \\
\bottomrule
\end{tabular}
\end{table}

\begin{finding}[Accelerator RF: Good Technical Fit, Niche Advantage]
\label{find:accel}
DFD P7 is a natural technical match for SRF accelerators: same cavity
technology, same cryogenic infrastructure, native RF coupling. The 98\%
round-trip efficiency at 20~ms exceeds conventional approaches.

The advantage is strongest for next-generation colliders requiring
$>$100~MJ pulse energy, where the high energy density reduces the
storage system footprint.

\textbf{Verdict}: \textbf{Good fit} for future SRF accelerators.
Limited by the fact that accelerator facilities already have SRF
infrastructure, so the incremental cost of P7 storage is lower than
greenfield deployment.
\end{finding}

\subsection{Application 5: Grid Frequency Regulation}

\subsubsection{Requirements}
\begin{itemize}[nosep]
  \item Response time: $<$200~ms (UK Grid Code), $<$500~ms (US PJM).
  \item Duration: 15--30~minutes of sustained output.
  \item Power: 1--50~MW per unit.
  \item Cost target: \$200--500/kW (competitive with gas peakers).
\end{itemize}

\subsubsection{DFD P7 performance}
\begin{itemize}[nosep]
  \item Response time: sub-$\mu$s (excellent, $100\times$ faster than
    required).
  \item Duration: $\tau = 1$~s passive. \textbf{Cannot sustain} 15--30
    minutes without active maintenance at $P_{\rm maintain} = E/\tau$.
  \item For 1~MW sustained output: $E = 1\ee{6}\times 1 = 1$~MJ
    maintained at $P_{\rm maintain} = 1$~MW. Net useful power = 0
    (all input goes to maintenance).
  \item The store can provide a \emph{single burst} of 46~MW for
    $\sim$1~s, but then is empty.
\end{itemize}

\begin{finding}[Grid Frequency Regulation: Not Viable]
\label{find:grid-freq}
Grid frequency regulation requires sustained output over 15--30 minutes.
DFD P7 with $\tau_{1/e} = 1$~s cannot sustain any net power output
beyond $\sim$3~s. The maintenance power equals the stored energy per
second, leaving no margin for useful output.

\textbf{Verdict}: \textbf{Eliminated}. Grid frequency regulation
requires flywheels ($\tau \sim$ minutes--hours), batteries ($\tau \sim$
hours), or SMES with much longer $\tau$ than P7.
\end{finding}

%=============================================================================
\section{Complete Application Space Map}
\label{sec:P7-map}
%=============================================================================

\begin{table}[h]
\centering
\caption{DFD P7 complete application space at $\tau_{1/e} = 1$~s,
$E_{\rm max} = 46$~MJ (baseline). Green = genuine advantage,
yellow = marginal, red = eliminated.}
\label{tab:P7-map}
\begin{tabular}{@{}p{3.5cm}cccl@{}}
\toprule
Application & $t_{\rm cycle}$ & $E$/pulse & DFD advantage? & Verdict \\
\midrule
\multicolumn{5}{@{}l}{\textbf{\textcolor{green!50!black}{GENUINE ADVANTAGE:}}} \\
Compact pulsed DEW & $\mu$s--ms & 1--50~MJ
  & Volume: $650\times$ & \textcolor{green!50!black}{High value} \\
IFE driver (10~Hz) & 0.1~s & 1--5~MJ
  & $\eta$: 99.9\%, rep rate & \textcolor{green!50!black}{Good fit} \\
Future accelerator RF & 1--20~ms & 10--100~MJ
  & Native SRF match & \textcolor{green!50!black}{Good fit} \\
\midrule
\multicolumn{5}{@{}l}{\textbf{\textcolor{orange!80!black}{MARGINAL:}}} \\
Railgun (baseline) & 2--5~ms & 32--64~MJ
  & Needs Nb$_3$Sn & \textcolor{orange!80!black}{Upgrade needed} \\
Regen braking buffer & 2--10~s & 1--5~MJ
  & Cryo impractical & \textcolor{orange!80!black}{Impractical} \\
EMP protection & $<$1~ms & 1--10~MJ
  & Niche & \textcolor{orange!80!black}{Possible} \\
\midrule
\multicolumn{5}{@{}l}{\textbf{\textcolor{red}{ELIMINATED:}}} \\
Grid storage (hrs) & $10^4$~s & GJ & $\eta \to 0$ & \textcolor{red}{Dead} \\
Grid freq regulation & 15--30~min & MW$\cdot$min & $\tau$ too short & \textcolor{red}{Dead} \\
UPS backup & 15~min & MJ & $\tau$ too short & \textcolor{red}{Dead} \\
EV range extension & $10^3$~s & MJ & $\tau$ too short & \textcolor{red}{Dead} \\
Solar overnight shift & $10^4$~s & GJ & $\eta \to 0$ & \textcolor{red}{Dead} \\
\bottomrule
\end{tabular}
\end{table}

%=============================================================================
\section{Head-to-Head Summary: P7 vs Competitors}
\label{sec:P7-head-to-head}
%=============================================================================

\begin{tcolorbox}[colback=blue!5!white, colframe=blue!60!black,
  title=\textbf{Where DFD P7 Wins}]
\begin{enumerate}[nosep]
\item \textbf{Energy density for pulsed discharge ($t < 100$~ms):}
  $650$~MJ/m$^3$ vs.\ $50$~MJ/m$^3$ (best flywheel) or $10$~MJ/m$^3$
  (best SMES). The $13$--$65{,}000\times$ density advantage is decisive
  when volume is the binding constraint.
\item \textbf{Round-trip efficiency at sub-ms timescales:}
  $>99.9\%$ vs.\ 90--95\% for SMES, 80--90\% for capacitor banks
  (including switching losses).
\item \textbf{Repetition rate at $>$10~MJ/pulse:}
  $\sim 4$~Hz at 10~MJ vs.\ $<0.1$~Hz for large capacitor banks
  (limited by recharge), $\sim 1$~Hz for SMES.
\end{enumerate}
\end{tcolorbox}

\begin{tcolorbox}[colback=red!5!white, colframe=red!60!black,
  title=\textbf{Where DFD P7 Loses}]
\begin{enumerate}[nosep]
\item \textbf{Any application requiring $t > 3$~s storage:}
  $\tau_{1/e} = 1$~s makes P7 useless for anything beyond burst
  applications. Flywheels ($\tau \sim$ hours), SMES ($\tau \sim$ hours
  with persistent-current mode), and batteries ($\tau \sim$ months)
  all dominate.
\item \textbf{System-level volume and mass:}
  The $\sim 110$~m$^3$ system footprint (including cryogenics) erases
  the cavity energy density advantage for applications where total
  system size matters more than cavity size.
\item \textbf{Cost and complexity:}
  SRF cavities at 2~K require He-II cryogenics. No commercial
  supply chain exists outside accelerator facilities. Capacitor banks
  and flywheels are commercial off-the-shelf.
\item \textbf{TRL:}
  P7 is at TRL~1. All competitors are TRL~5--9.
\end{enumerate}
\end{tcolorbox}

%=============================================================================
\part{Patent 12: Break-Even Economics for City-Scale Psi-Transmission}
\label{part:P12}
%=============================================================================

%=============================================================================
\section{W3i6 Network Parameters}
\label{sec:P12-params}
%=============================================================================

From W3i6 P4+P12 update:
\begin{center}
\begin{tabular}{ll}
\toprule
Parameter & W3i6 Value \\
\midrule
End-to-end efficiency $\eta_{\rm e2e}$ & 20\% (= 50\% hub $\times$ 81\%
  transport $\times$ 50\% endpoint, rounded) \\
Capital cost $C_{\rm cap}$ & \$0.6--1.1~B \\
Network topology & Binary tree, 50~km radius \\
Hub: SRF (Nb$_3$Sn, 4~K) & \$80~M \\
Corridor (passive) & \$20~M \\
Endpoints ($\times 100$, SRF) & \$500~M--\$1~B \\
Transport efficiency (corridor only) & 81\% \\
Conversion efficiency (each stage) & 50\% (at polariton crossing, SRF) \\
\bottomrule
\end{tabular}
\end{center}

%=============================================================================
\section{Levelised Cost Model}
\label{sec:P12-LCOE}
%=============================================================================

\subsection{Framework}

The levelised cost of energy delivery (LCOD) for the psi-network includes:
\begin{enumerate}[nosep]
  \item Capital cost amortised over lifetime $T$.
  \item Energy losses (input energy exceeds delivered energy by $1/\eta_{\rm e2e}$).
  \item Operations and maintenance (O\&M).
\end{enumerate}

The LCOD in \$/kWh delivered:
\begin{equation}
\text{LCOD}_\psi = \frac{C_{\rm cap}\,\text{CRF}}{P_{\rm del}\times 8760\times\text{CF}}
+ \frac{p_{\rm elec}}{\eta_{\rm e2e}}
+ \text{O\&M}_{\rm fixed},
\label{eq:LCOD}
\end{equation}
where:
\begin{itemize}[nosep]
  \item CRF $= r(1+r)^T/[(1+r)^T - 1]$ is the capital recovery factor.
  \item $P_{\rm del}$ is the rated delivered power (MW).
  \item CF is the capacity factor (fraction of time at rated power).
  \item $p_{\rm elec}$ is the input electricity price (\$/kWh).
  \item $\eta_{\rm e2e} = 0.20$ is the end-to-end efficiency.
  \item O\&M$_{\rm fixed}$ covers cryogenic maintenance, SRF cavity
    refurbishment, and staffing.
\end{itemize}

\subsection{Reference parameters}

\begin{center}
\begin{tabular}{lll}
\toprule
Parameter & Value & Basis \\
\midrule
$C_{\rm cap}$ & \$0.6~B (low) / \$1.1~B (high) & W3i6 \\
$T$ (lifetime) & 30~years & Comparable to grid infrastructure \\
$r$ (discount rate) & 8\% & Utility-scale infrastructure \\
CRF & 0.0888/yr & From $r$ and $T$ \\
$P_{\rm del}$ & 10~GW & City-scale (London/NYC class) \\
CF & 0.60 & Average grid capacity factor \\
$p_{\rm elec}$ (input) & \$0.05/kWh & Wholesale \\
$\eta_{\rm e2e}$ & 0.20 & W3i6 \\
O\&M & 2\% of $C_{\rm cap}$/yr & Standard for major infrastructure \\
\bottomrule
\end{tabular}
\end{center}

\subsection{LCOD calculation}

\subsubsection{Capital component}
\begin{equation}
\text{LCOD}_{\rm cap} = \frac{C_{\rm cap}\times\text{CRF}}{P_{\rm del}
\times 8760\times\text{CF}}.
\end{equation}

At $C_{\rm cap} = \$0.6$~B:
\begin{equation}
\text{LCOD}_{\rm cap}^{\rm low} = \frac{6\ee{8}\times 0.0888}
{10\ee{6}\times 8760\times 0.60}
= \frac{5.33\ee{7}}{5.26\ee{10}}
= 0.00101~\$/\text{kWh}
= 0.10~\text{c/kWh}.
\end{equation}

At $C_{\rm cap} = \$1.1$~B:
\begin{equation}
\text{LCOD}_{\rm cap}^{\rm high} = \frac{1.1\ee{9}\times 0.0888}
{5.26\ee{10}}
= \frac{9.77\ee{7}}{5.26\ee{10}}
= 0.00186~\$/\text{kWh}
= 0.19~\text{c/kWh}.
\end{equation}

\textbf{Note}: The capital component is surprisingly small ($<$0.2~c/kWh)
because the 10~GW throughput amortises the cost over an enormous energy
volume ($5.26\ee{10}$~kWh/yr).

\subsubsection{Energy loss component}

This is the dominant cost. The network delivers only 20\% of input energy:
\begin{equation}
\text{LCOD}_{\rm loss} = \frac{p_{\rm elec}}{\eta_{\rm e2e}}
= \frac{0.05}{0.20} = 0.25~\$/\text{kWh}.
\end{equation}

The psi-network must \emph{buy} 5~kWh of input electricity for every
1~kWh delivered. At \$0.05/kWh wholesale, the energy cost alone is
\$0.25/kWh.

\subsubsection{O\&M component}
\begin{equation}
\text{O\&M} = \frac{0.02\times C_{\rm cap}}{P_{\rm del}\times 8760\times\text{CF}}.
\end{equation}

At $C_{\rm cap} = \$0.8$~B (midpoint):
\begin{equation}
\text{O\&M} = \frac{0.02\times 8\ee{8}}{5.26\ee{10}}
= 0.00030~\$/\text{kWh} = 0.03~\text{c/kWh}.
\end{equation}

Negligible.

\subsubsection{Total LCOD}

\begin{equation}
\boxed{
\text{LCOD}_\psi = \underbrace{0.10\text{--}0.19}_{\rm capital}
+ \underbrace{25.0}_{\rm energy~loss}
+ \underbrace{0.03}_{\rm O\&M}
= 25.1\text{--}25.2~\text{c/kWh}}
\quad\text{at $p_{\rm elec} = 5$~c/kWh.}
\end{equation}

\begin{finding}[LCOD of Psi-Network at 20\% Efficiency]
\label{find:LCOD}
The levelised cost of energy delivery for the W3i6 city-scale psi-network
is:
\begin{equation}
\text{LCOD}_\psi \approx \frac{p_{\rm elec}}{0.20} + 0.001\text{--}0.002
\approx 5\,p_{\rm elec} + \epsilon,
\end{equation}
where $\epsilon \sim 0.2$~c/kWh covers capital and O\&M (negligible at
high throughput).

\textbf{The LCOD is dominated by energy losses}: the $5\times$ input
energy multiplier from 20\% efficiency. Capital cost is $<$1\% of total
LCOD at city-scale throughput.
\end{finding}

%=============================================================================
\section{Break-Even Analysis}
\label{sec:P12-breakeven}
%=============================================================================

\subsection{Break-even condition}

The psi-network becomes viable when its LCOD is less than the alternative
cost of conventional electricity delivery:
\begin{equation}
\text{LCOD}_\psi < p_{\rm conv},
\end{equation}
where $p_{\rm conv}$ is the all-in cost of conventional grid delivery
to the same endpoints (\$/kWh).

\begin{equation}
\frac{p_{\rm elec}}{\eta_{\rm e2e}} + \frac{C_{\rm cap}\times\text{CRF}}
{P_{\rm del}\times 8760\times\text{CF}} < p_{\rm conv}.
\label{eq:breakeven}
\end{equation}

Since the capital term is negligible ($\sim$0.2~c/kWh):
\begin{equation}
\boxed{
p_{\rm conv}^{\rm break\text{-}even} \approx
\frac{p_{\rm elec}}{\eta_{\rm e2e}} = 5\,p_{\rm elec}.
}
\label{eq:breakeven-simple}
\end{equation}

\subsection{Break-even electricity prices}

\begin{table}[h]
\centering
\caption{Break-even conventional delivery price for psi-network viability,
at various input electricity costs.}
\label{tab:breakeven}
\begin{tabular}{ccc}
\toprule
Input electricity $p_{\rm elec}$ & Required $p_{\rm conv}^{\rm break\text{-}even}$
  & Context \\
 (c/kWh) & (c/kWh) & \\
\midrule
2 (ultra-cheap hydro/nuclear) & 10 & $= $ US average retail \\
3 (cheap wholesale) & 15 & 1.5$\times$ US average \\
5 (US wholesale average) & 25 & 2.5$\times$ US average \\
8 (EU wholesale) & 40 & Most expensive island grids \\
10 (premium wholesale) & 50 & Exceeds all current markets \\
\bottomrule
\end{tabular}
\end{table}

\begin{finding}[Break-Even Conventional Price]
\label{find:breakeven}
The psi-network breaks even when conventional delivery costs exceed
$5\times$ the wholesale electricity input price:

\begin{itemize}[nosep]
  \item At US wholesale input (\$0.05/kWh): break-even at $p_{\rm conv}
    > \$0.25$/kWh. This is $2.5\times$ the US retail average (\$0.10/kWh)
    but within the range of some island nations (Hawaii \$0.28, US Virgin
    Islands \$0.35).
  \item At ultra-cheap input (\$0.02/kWh, e.g., stranded hydro):
    break-even at $p_{\rm conv} > \$0.10$/kWh --- exactly the US average
    retail price.
  \item At EU wholesale (\$0.08/kWh): break-even at $p_{\rm conv}
    > \$0.40$/kWh. No current market has such high delivery costs.
\end{itemize}

\textbf{The narrow viability window}: psi-transmission is viable only
when cheap input electricity is available AND the delivery endpoint has
very expensive conventional alternatives.
\end{finding}

%=============================================================================
\section{Sensitivity Analysis}
\label{sec:P12-sensitivity}
%=============================================================================

\subsection{Sensitivity to efficiency}

The LCOD is inversely proportional to $\eta_{\rm e2e}$ (for the
dominant energy-loss component):
\begin{equation}
\text{LCOD}_\psi \approx \frac{p_{\rm elec}}{\eta_{\rm e2e}}.
\end{equation}

\begin{table}[h]
\centering
\caption{LCOD sensitivity to efficiency improvement, at $p_{\rm elec}
= \$0.05$/kWh.}
\label{tab:sensitivity-eta}
\begin{tabular}{ccl}
\toprule
$\eta_{\rm e2e}$ & LCOD (c/kWh) & How to achieve \\
\midrule
20\% (W3i6 baseline) & 25.0 & Two 50\% conversion stages \\
33\% & 15.2 & One 50\% conversion + 65\% transport \\
50\% & 10.0 & Direct SRF-to-SRF transfer (no polariton) \\
65\% & 7.7 & Single-stage 80\% conversion \\
81\% & 6.2 & Zero conversion loss (transport only) \\
\bottomrule
\end{tabular}
\end{table}

\begin{finding}[Efficiency is the Critical Lever]
\label{find:efficiency-lever}
Every percentage point of efficiency improvement reduces LCOD by
$\sim p_{\rm elec}/\eta^2$. The most impactful improvement is
\textbf{eliminating one conversion stage}:

\begin{itemize}[nosep]
  \item If the hub conversion could be eliminated (e.g., if the
    generation source directly produces $\dpsiEM$ --- a ``psi-mode
    generator''), efficiency jumps from 20\% to $\sim$40\%, halving LCOD.
  \item If \emph{both} conversions are eliminated (psi-mode generation
    at source, psi-mode consumption at load), efficiency reaches 81\%
    (transport only), and LCOD drops to 6.2~c/kWh --- competitive with
    conventional grid at any input price.
\end{itemize}

The physics of direct psi-mode generation (bypassing EM$\to\dpsiEM$
conversion) is an open question that merits investigation.
\end{finding}

\subsection{Sensitivity to throughput}

\begin{table}[h]
\centering
\caption{LCOD sensitivity to network throughput, at $C_{\rm cap} = \$0.8$~B,
$\eta = 20\%$, $p_{\rm elec} = \$0.05$/kWh.}
\label{tab:sensitivity-throughput}
\begin{tabular}{ccc}
\toprule
$P_{\rm del}$ (GW) & Capital LCOD (c/kWh) & Total LCOD (c/kWh) \\
\midrule
0.1 & 13.5 & 38.5 \\
1.0 & 1.35 & 26.4 \\
10 & 0.14 & 25.1 \\
100 & 0.01 & 25.0 \\
\bottomrule
\end{tabular}
\end{table}

Capital cost matters only below $\sim$1~GW throughput. At city scale
($\geq$10~GW), the economics are entirely determined by the efficiency.

\subsection{Sensitivity to capital cost}

\begin{table}[h]
\centering
\caption{LCOD sensitivity to capital cost, at 10~GW throughput.}
\label{tab:sensitivity-cap}
\begin{tabular}{ccc}
\toprule
$C_{\rm cap}$ (\$B) & Capital LCOD (c/kWh) & Total LCOD (c/kWh) \\
\midrule
0.1 & 0.02 & 25.0 \\
0.6 & 0.10 & 25.1 \\
1.1 & 0.19 & 25.2 \\
5.0 & 0.85 & 25.9 \\
\bottomrule
\end{tabular}
\end{table}

\begin{finding}[Capital Cost is Not the Binding Constraint]
\label{find:capital-not-binding}
At city-scale throughput ($\geq$10~GW), even a $5\times$ increase in
capital cost (to \$5~B) adds only 0.85~c/kWh to LCOD. A $5\times$
decrease (to \$0.2~B) saves only 0.08~c/kWh. The capital cost is
irrelevant to economic viability.

\textbf{The binding constraint is efficiency}. No amount of capital
cost reduction makes the network viable if $\eta_{\rm e2e}$ remains
at 20\%.
\end{finding}

%=============================================================================
\section{Comparison with Underground WPT (P4)}
\label{sec:P12-vs-P4}
%=============================================================================

The P4 underground WPT system provides a useful benchmark:

\begin{table}[h]
\centering
\caption{P4 underground WPT vs.\ P12 city-scale network.}
\label{tab:P4-vs-P12}
\begin{tabular}{@{}lcc@{}}
\toprule
Parameter & P4 (Underground WPT) & P12 (City Network) \\
\midrule
Efficiency & 75.2\% & 20\% \\
Capital & \$270~M & \$600~M--1,100~M \\
LCOD (at 5~c/kWh input) & 6.6~c/kWh & 25~c/kWh \\
Conversion stages & 2 (but 50\% each $\to$ 25\%) & 2 (50\% each $\to$ 25\%) \\
Transport efficiency & $\approx 100\%$ & 81\% \\
Topology & Point-to-point & Distributed tree \\
Break-even $p_{\rm conv}$ & 6.6~c/kWh & 25~c/kWh \\
\bottomrule
\end{tabular}
\end{table}

Wait --- the P4 efficiency of 75.2\% is the \emph{transport} efficiency
(beam capture). It also requires two conversion stages (EM$\to\dpsiEM$
at transmitter, $\dpsiEM\to$EM at receiver), each at 50\%. So the true
P4 end-to-end efficiency is:
\begin{equation}
\eta_{\rm P4}^{\rm e2e} = 50\% \times 75.2\% \times 50\% = 18.8\%.
\end{equation}

This is essentially identical to the P12 efficiency of 20\% (the
difference is $81\%$ vs.\ $75.2\%$ transport). The conversion
bottleneck affects \emph{both} patents equally.

\begin{finding}[P4 and P12 Share the Same Conversion Bottleneck]
\label{find:shared-bottleneck}
The end-to-end efficiency of both P4 (18.8\%) and P12 (20\%) is
dominated by the two 50\% conversion stages. The transport stage
(75.2\% for P4 beam capture, 81\% for P12 corridor) is relatively
efficient in both cases.

\textbf{Implication}: The W3i6 P4 LCOE claim of \$0.004/kWh assumed
$\eta = 75.2\%$ for the full system, but this is the transport-only
efficiency. The corrected P4 end-to-end LCOE at 18.8\% is:
\begin{equation}
\text{LCOE}_{\rm P4}^{\rm corrected} = \frac{p_{\rm elec}}{\eta_{\rm e2e}}
+ \frac{C_{\rm cap}\times\text{CRF}}{P_{\rm del}\times 8760\times\text{CF}}
\approx \frac{0.05}{0.188} + \epsilon = 26.6~\text{c/kWh}.
\end{equation}

\textbf{This is a significant correction to the P4 economics.} The
\$0.004/kWh figure assumed that the 75.2\% transport efficiency was the
total efficiency. With conversion losses included, P4 and P12 have
essentially the same LCOD.
\end{finding}

\begin{tcolorbox}[colback=red!5!white, colframe=red!60!black,
  title=\textbf{CORRECTION: P4 LCOE Requires Conversion Losses}]
The W3i6 P4 LCOE of \$0.004/kWh \textbf{omitted the two conversion
stages} (EM$\to\dpsiEM$ at transmitter, $\dpsiEM\to$EM at receiver).
Including these at 50\% each:
\begin{equation}
\eta_{\rm P4}^{\rm e2e} = 0.50\times 0.752\times 0.50 = 18.8\%.
\end{equation}
The corrected P4 LCOE is $\sim$\$0.27/kWh at \$0.05/kWh input ---
\textbf{comparable to the P12 city-scale network}.

Both P4 and P12 are economically limited by the same conversion
bottleneck. Breaking this bottleneck (through higher $\lambdaone$,
parametric amplification, or direct psi-mode generation) is the
single most important technology milestone for all DFD energy
transmission patents.
\end{tcolorbox}

%=============================================================================
\section{Niche Markets: Where Psi-Transmission Could Be Viable}
\label{sec:P12-niche}
%=============================================================================

Despite the unfavourable economics at 20\% efficiency, certain markets
could justify the cost:

\subsection{Market 1: Undersea power delivery}

\begin{itemize}[nosep]
  \item Deep-sea mining, ocean-floor observatories, military
    installations require power at locations where conventional cables
    are extremely expensive (\$1--10~M/km for subsea HVDC).
  \item A 100~km subsea cable at \$5~M/km = \$500~M capital.
    Psi-corridor: no physical cable, \$80~M (single SRF hub + endpoint).
  \item At 20\% efficiency and \$0.05/kWh input, LCOD = 25~c/kWh.
    But the alternative (diesel generator on-site) costs \$0.30--1.00/kWh.
  \item \textbf{Viable} for remote subsea applications where the
    alternative is on-site generation at $>$\$0.30/kWh.
\end{itemize}

\subsection{Market 2: Deep underground mining}

\begin{itemize}[nosep]
  \item Ultra-deep mines (3--5~km depth, South African gold mines)
    spend \$0.15--0.25/kWh on ventilation, cooling, and power delivery
    infrastructure.
  \item Psi-transmission bypasses all physical conduit requirements.
  \item Break-even: at ultra-cheap surface power (\$0.02/kWh),
    LCOD = 10~c/kWh. Competitive if underground delivery infrastructure
    costs $>$10~c/kWh.
  \item \textbf{Marginally viable} for extreme-depth mining.
\end{itemize}

\subsection{Market 3: Space-based power delivery}

\begin{itemize}[nosep]
  \item Lunar or Martian surface bases: no grid infrastructure.
    Alternative is nuclear or solar with battery storage.
  \item Psi-corridor from orbital solar power station to surface:
    no atmospheric losses, no physical infrastructure.
  \item Cost comparison: space solar panel \$500--5000/W, nuclear
    \$100--500/W. Psi-transmission capital is amortised differently
    in space contexts.
  \item \textbf{Potentially viable} in long-term space infrastructure
    scenarios (post-2050).
\end{itemize}

\begin{finding}[Niche Market Viability]
\label{find:niche}
The city-scale psi-network is not viable for general urban power
delivery at current efficiencies. However, niche markets exist where:
\begin{enumerate}[nosep]
  \item Conventional delivery is physically impossible or extremely
    expensive ($>$\$0.25/kWh delivered).
  \item The psi-corridor's unique capability (no physical conduit,
    underground/undersea penetration) provides irreplaceable value.
  \item Input power is available at ultra-low cost (\$0.02--0.03/kWh).
\end{enumerate}

\textbf{Best niche candidates}: undersea power delivery, deep mining,
and long-term space infrastructure.
\end{finding}

%=============================================================================
\section{Technology Milestones for Economic Viability}
\label{sec:P12-milestones}
%=============================================================================

\begin{table}[h]
\centering
\caption{Technology milestones and their impact on psi-network LCOD.
All at $p_{\rm elec} = \$0.05$/kWh, 10~GW throughput.}
\label{tab:milestones}
\begin{tabular}{@{}p{5cm}ccc@{}}
\toprule
Milestone & $\eta_{\rm e2e}$ & LCOD (c/kWh) & Viable? \\
\midrule
W3i6 baseline (SRF conversion) & 20\% & 25.0 & \textcolor{red}{No} \\
Improved SRF ($Q > 10^{11}$) & 30\% & 16.7 & \textcolor{red}{No} \\
Single-stage conversion & 40\% & 12.5 & \textcolor{orange}{Marginal} \\
High-$\lambdaone$ material & 50\% & 10.0 & \textcolor{orange}{Marginal} \\
Direct psi-mode generation & 65\% & 7.7 & \textcolor{green!50!black}{Competitive} \\
Lossless conversion & 81\% & 6.2 & \textcolor{green!50!black}{Yes} \\
\bottomrule
\end{tabular}
\end{table}

\begin{finding}[Path to Economic Viability]
\label{find:viability-path}
The psi-network becomes economically competitive with conventional
urban grids (\$0.06--0.12/kWh) only when $\eta_{\rm e2e} \geq 65\%$,
which requires:
\begin{enumerate}[nosep]
  \item Elimination of at least one conversion stage (direct psi-mode
    generation or direct psi-mode consumption), AND
  \item Remaining conversion stage at $\geq 80\%$ efficiency.
\end{enumerate}

Alternatively, at the current 20\% efficiency, break-even requires
conventional delivery costs $>$25~c/kWh --- achievable only in extreme
niche markets.

\textbf{The conversion efficiency is the gating technology for all DFD
energy transmission economics.} This is the single most important
research target for the DFD energy patent portfolio.
\end{finding}

%=============================================================================
\part{Cross-Patent Integration and Open Questions}
\label{part:cross}
%=============================================================================

%=============================================================================
\section{Cross-References}
\label{sec:crossrefs}
%=============================================================================

\subsection{P4 $\leftrightarrow$ P12: Shared conversion bottleneck}

The W3i7 analysis reveals that P4 and P12 are economically limited by
the same conversion stages. The W3i6 P4 LCOE of \$0.004/kWh is
corrected to \$0.27/kWh (at \$0.05 input). Both patents require the
same technology breakthrough: improved EM-to-$\dpsiEM$ conversion.

\subsection{P4 $\leftrightarrow$ P2: Parametric amplification dead end}

The W3i6 open question about parametric pumping of $\dpsiEM$ via
$\psigrav$ gradient is now answered: the mechanism is real but produces
negligible gain at terrestrial gravity ($\gamma_{\rm par} \sim 10^{-14}$
m$^{-1}$). The P2 parametric amplification patent would need to operate
through a different mechanism (e.g., temporal modulation of the cavity
boundary conditions, or EM-pump-driven MOND nonlinearity) rather than
static gravitational gradient pumping.

\subsection{P7 $\leftrightarrow$ P11: SRF technology synergy}

P7 and P11 share the SRF cavity technology base. The P7 application
space (pulsed power buffer) could leverage P11 SRF cavity production.
The Nb$_3$Sn upgrade path (W3i4 P11) would increase P7 $E_{\rm max}$
from 46~MJ to $\sim$208~MJ, significantly expanding the railgun and
accelerator applications.

\subsection{P7 $\leftrightarrow$ P4/P12: No synergy}

P7 (energy storage) and P4/P12 (energy transmission) address different
problems. P7 cannot serve as a buffer in the P12 network because the
1~s storage time is insufficient for network latency management
(typical grid contingency requires 15~min reserves).

%=============================================================================
\section{Open Questions for W3i8}
\label{sec:open}
%=============================================================================

\begin{enumerate}
\item \textbf{P4/P12: Direct psi-mode generation.} Is there a physical
  mechanism to generate $\dpsiEM$ directly (without EM$\to\dpsiEM$
  conversion)? This would eliminate the conversion bottleneck. Candidates:
  piezoelectric-analogue coupling in MOND materials, direct mechanical
  excitation of the psi-mode in SRF cavities, or stimulated psi-mode
  emission.

\item \textbf{P4/P12: Conversion efficiency correction propagation.}
  The P4 LCOE of \$0.004/kWh is now identified as transport-only.
  Propagate this correction to all P4 economic claims and the P4 patent
  specification.

\item \textbf{P7: Nb$_3$Sn upgrade quantification.} With $B_{\rm sh}
  = 425$~mT (Nb$_3$Sn), calculate the revised $\Pmax$, $E_{\rm max}$,
  and updated application competitiveness for railgun and CLIC-class
  accelerator.

\item \textbf{P7: Cryogenic packaging for DEW.} The DEW application
  hinges on miniaturising He-II cryogenics to $<$5~m$^3$. What is the
  state of the art in compact 2~K cryocoolers? Can pulse-tube systems
  reach 2~K in the required volume?

\item \textbf{P12: Undersea market sizing.} Quantify the addressable
  market for undersea psi-power delivery: deep-sea mining installations,
  ocean-floor observatories, military assets. What is the annual energy
  demand and willingness-to-pay?

\item \textbf{P2 $\leftrightarrow$ P4: Alternative parametric mechanisms.}
  Since static $\psigrav$-gradient pumping is a dead end, investigate
  temporal modulation of cavity boundary conditions (dynamic Casimir-type
  effect in the psi-sector) as an alternative parametric amplification
  pathway.
\end{enumerate}

%=============================================================================
\section{Conclusion}
\label{sec:conclusion}
%=============================================================================

\begin{tcolorbox}[colback=green!5!white, colframe=green!60!black,
  title=\textbf{W3i7 P4+P7+P12 Bottom Line}]

\textbf{P4 (Parametric Pumping):}
\begin{itemize}[nosep]
  \item Parametric gain formula derived:
    $G = \cosh^2[(1+\kappa)g\omega L/(c^3\npsi)]$.
  \item At Earth's surface: $\gamma_{\rm par} = 3.35\ee{-14}$~m$^{-1}$
    at 2.45~GHz. Gain is negligible ($G - 1 \sim 10^{-27}$ per metre).
  \item \textbf{Cannot circumvent SRF conversion bottleneck.}
    Required interaction length: 5.8~AU. Even with resonant enhancement
    ($\Qpsi = 10^9$): $G - 1 \sim 10^{-13}$.
  \item Useful only at neutron-star gravity ($g \sim 10^{12}$~m/s$^2$).
  \item P4 underground WPT specification unchanged from W3i6.
\end{itemize}

\medskip

\textbf{P7 (Complete Application Space):}
\begin{itemize}[nosep]
  \item \textbf{Genuine advantages}: compact DEW/HPM (volume-decisive),
    IFE driver (99.9\% $\eta$, 10~Hz), future accelerator RF fill.
  \item \textbf{Marginal}: railgun (needs Nb$_3$Sn upgrade), EMP
    protection.
  \item \textbf{Eliminated}: grid storage, grid frequency regulation,
    UPS, EV, solar shifting --- all require $\tau \gg 1$~s.
  \item The $650$~MJ/m$^3$ density advantage is real but only
    translates to system-level advantage when cryogenic packaging is
    compact ($<$5~m$^3$).
\end{itemize}

\medskip

\textbf{P12 (Break-Even Economics):}
\begin{itemize}[nosep]
  \item LCOD $\approx 5\times p_{\rm elec}$ at 20\% efficiency.
    At \$0.05/kWh input: LCOD = 25~c/kWh.
  \item Break-even: $p_{\rm conv} > 25$~c/kWh (niche markets only).
  \item Capital cost is $<$1\% of LCOD at city scale. \textbf{Efficiency
    is the only lever that matters.}
  \item Viable niches: undersea, deep mining, space.
  \item \textbf{CORRECTION}: P4 LCOE of \$0.004/kWh omitted conversion
    losses. Corrected P4 LCOE $\approx$ 27~c/kWh. P4 and P12 share
    the same conversion bottleneck.
  \item Path to competitiveness: $\eta_{\rm e2e} \geq 65\%$ (requires
    eliminating at least one conversion stage).
\end{itemize}

\medskip

\textbf{Overarching finding}: The EM-to-$\dpsiEM$ conversion efficiency
is the \textbf{single gating technology} for all DFD energy patents
(P4, P7, P12). Parametric pumping via $\psigrav$ is not a viable
workaround. The research priority must be: (a) improved SRF $Q$ for
higher resonant buildup, (b) discovery of higher-$\lambdaone$ materials,
or (c) direct psi-mode generation bypassing EM conversion entirely.
\end{tcolorbox}

\end{document}
